\documentclass{beamer}

% Using your requested theme and color
\usetheme{Szeged}
\usecolortheme{beaver}

% --- Presentation Info ---
\title{Fibonacci Heap}
\subtitle{Outperforming Traditional Priority Queues}
\author{Pedro Marinho \and Miguel Mateus \and Tuomas Haapasalo}
\institute{FEUP}
\date{March 23, 2026}

% Removes the small navigation icons at the bottom right
\setbeamertemplate{navigation symbols}{}
\setbeamertemplate{itemize items}[triangle]
\setbeamercolor{itemize item}{fg=black}
\setbeamercolor{itemize subitem}{fg=black}
\begin{document}

% --- Title Slide ---
\begin{frame}
  \titlepage
\end{frame}

% --- 1. OUTLINE (1 dot) ---
\section{Overview}
\begin{frame}{Overview}
  % sections={2-8} ensures the "Outline" section doesn't list itself
  \tableofcontents[sections={2-8}]
\end{frame}

% --- AUTOMATED TOPIC TRANSITIONS ---
% Placing this here means it will automatically generate the grayed-out 
% transition slide for every section AFTER the Outline.
\AtBeginSection[]
{
  \begin{frame}{Overview}
    % currentsection highlights the active topic and fades the rest
    \tableofcontents[currentsection, sections={2-8}]
  \end{frame}
}

% --- 2. MOTIVATION (1 dot) ---
\section{Motivation}

\begin{frame}{Motivation}
  \begin{itemize}
    \item Efficient priority queues are a key problem in the context of graph algorithm design.
    \item There is an increased need for processing massive networks, but the repeated cost of updating node priorities poses a huge bottleneck!
    \item Demand for priority queue designs capable of extremely fast updates, but that can simultaneously keep other standard operations efficient.
  \end{itemize}
  
  \vspace{0.2cm}
  {\large \textcolor{darkred}{The $O(\log n)$ Bound}}
  
  \vspace{0.1cm}
  Currently, the effects of logarithmic slowness are much more visible! The number of edges in dense networks vastly exceeds the vertices, making priority updates the dominant computing cost \cite{fredman1987}.
\end{frame}

\section{Data Structure Basis}
\begin{frame}[c]{The Fibonacci Heap}{Definition}
  
  \vspace{0.4cm}

  The Fibonacci heap is a data structure consisting of a collection,
  or forest, of min-heap ordered trees, such that:
  
  \vspace{0.1cm}
  
  \begin{itemize}
    \item Every parent node maintains a key value that is less than or equal to the values of its children
    \item The absolute minimum element of any individual tree will always be found at its root
    \item The trees in this forest do not need to maintain a rigid or balanced shape
  \end{itemize}
  
  \vspace{0.4cm}
  
  \begin{center}
    
  \end{center}

\end{frame}

\begin{frame}[c]{The Fibonacci Heap}{Anatomy of a Node}

  \vspace{0.4cm}

  Each node $x$ in the heap contains several important attributes:
  
  \vspace{0.1cm}
  
  \begin{itemize}
    \item $x.key$: Priority value of the element
    \item $x.degree$: Number of direct children the node possesses
    \item Pointers:
      \begin{itemize}
        \item Vertical: $x.parent$ and $x.child$ (points to any one of its children)
        \item Horizontal: $x.left$ and $x.right$ (connects to adjacent siblings)
      \end{itemize}
    \item $x.mark$: A boolean value indicating whether the node has lost a child since it was last made the child of another node
  \end{itemize}

  \vspace{0.2cm}
  \begin{center}
    
  \end{center}

\end{frame}


\begin{frame}[c]{The Fibonacci Heap}{Circular Doubly Linked Lists}

  \vspace{0.4cm}
  
  To organize the nodes efficiently, this data structure relies on circular doubly linked lists
  
  \vspace{0.1cm}
  
  \begin{itemize}
    \item All siblings, which are children of the same parent, are connected to each other in a continuous loop
    \item The roots of all the trees in the forest are also linked together to form a global root list
  \end{itemize}
  
  \vspace{0.4cm}
  
  {\large \textcolor{darkred}{Why this matters}}
  
  \vspace{0.1cm}
  
  This way of organizing allows us to more easily manage and compute operations, such as insertion or key decrease, in an amortized time of $O(1)$.
  
  \vspace{0.2cm}
  
  \begin{center}
    
  \end{center}

\end{frame}


\begin{frame}[c]{The Fibonacci Heap}{Tracking the Minimum}

  \vspace{0.6cm}
  
  To quickly access the minimum element without searching the structure:
  
  \vspace{0.1cm}
  
  \begin{itemize}
    \item A single global pointer, referred to as $min$, is maintained at all times
    \item This pointer always directs to the root node containing the absolute minimum key in the entire forest
    \item This direct access guarantees that finding the minimum element takes a worst-case time of $O(1)$
  \end{itemize}
  
  \vspace{0.4cm}
  
  \begin{center}
    
  \end{center}

\end{frame}


\begin{frame}[c]{The Lazy Philosophy}

  \vspace{0.6cm}
  
  The core design of this data structure relies heavily on intentionally delaying computational work
  
  \vspace{0.1cm}
  
  \begin{itemize}
    \item When inserting nodes or decreasing keys, the heap entirely avoids immediate restructuring
    \item New or modified nodes are simply added to the root list in constant $O(1)$ time
    \item Structural cleanup and tree merging are delayed until the extract-minimum operation is specifically called
  \end{itemize}
  
  \vspace{0.4cm}
  
  \begin{center}
    
  \end{center}

\end{frame}

% --- 4. OPERATIONS (9 dots) ---
\section{Operations}
\begin{frame}{Operations - Slide 1}\end{frame}
\begin{frame}{Operations - Slide 2}\end{frame}
\begin{frame}{Operations - Slide 3}\end{frame}
\begin{frame}{Operations - Slide 4}\end{frame}
\begin{frame}{Operations - Slide 5}\end{frame}


% --- 5. APPLICATIONS (4 dots) ---
\section{Applications}
\begin{frame}{Applications - Slide 1}\end{frame}
\begin{frame}{Applications - Slide 2}\end{frame}
\begin{frame}{Applications - Slide 3}\end{frame}
\begin{frame}{Applications - Slide 4}\end{frame}

% --- 6. EXAMPLE (4 dots) ---
\section{Example}
\begin{frame}{Example - Slide 1}\end{frame}
\begin{frame}{Example - Slide 2}\end{frame}
\begin{frame}{Example - Slide 3}\end{frame}
\begin{frame}{Example - Slide 4}\end{frame}

% --- 7. CONCLUSIONS (2 dots) ---
\section{Conclusions}
\begin{frame}{Conclusions - Slide 1}\end{frame}
\begin{frame}{Conclusions - Slide 2}\end{frame}

% --- 8. REFERENCES (2 dots) ---
\section{References}
\begin{frame}{References}
  \begin{thebibliography}{99}
    
    \bibitem{fredman1987}
    Michael L. Fredman and Robert Endre Tarjan.
    \newblock Fibonacci heaps and their uses in improved network optimization algorithms.
    \newblock \emph{Journal of the ACM}, 34(3):596--615, 1987.
    
  \end{thebibliography}
\end{frame}
\end{document}